\section[Categorías fibra]{Categorías fibra; equivalencia de
$\mathcal{E}$-categorías}
\label{VI.4}

Sea~$\cal{F}$ una categoría sobre~$\cal{E}$, y sea $S\in
\Ob({\cal{E}})$. Se llama \emph{categoría fibra de}~$\cal{F}$
en~$S$
\index{categoría fibra|hyperpage}%
a la subcategoría~$\cal{F}_S$
\label{indnot:fn}\oldindexnot{$\cal{F}_S$|hyperpage}%
de~$\cal{F}$ que es imagen inversa de la
\marginpar{159}
subcategoría puntual de~$\cal{E}$ definida por~$S$. Por tanto, los objetos
de~$\cal{F}_S$ son los objetos~$\xi$ de~$\cal{F}$ tales que $p(\xi)=S$, y
sus morfismos son los morfismos~$u$ de~$\cal{F}$ tales que
$p(u)=\id_S$, es decir, los $S$-morfismos en~$\cal{F}$. Naturalmente,
$\cal{F}_S$ es canónicamente isomorfa al producto fibrado
$\cal{F}\times_{\cal{E}}\{S\}$, donde~$\{S\}$ designa la subcategoría
puntual de~$\cal{E}$ definida por~$S$, provista de su funtor de inclusión
en~$\cal{E}$. Se deduce de ello (habida cuenta de la transitividad del cambio
de base) que, si se efectúa el cambio de base
$\lambda\colon\cal{E}'\to\cal{E}$, entonces, para todo
$S'\in\Ob(\cal{E}')$, \emph{la proyección $\pr_1\colon
\cal{F}'=\cal{F}\times_{\cal{E}}\cal{E}'\to\cal{F}$ induce un
isomorfismo}
$$
\cal{F}'_{S'}\to\cal{F}_S\qquad(\text{donde $S=\lambda(S')$})%
\ifthenelse{\boolean{orig}}{}{.}
$$

\begin{proposition}
\label{VI.4.1} Sea $f\colon\cal{F}\to\cal{G}$ un
$\cal{E}$-funtor. Si~$f$ es plenamente fiel, entonces, para todo cambio de
base $\cal{E}'\to\cal{E}$, el funtor correspondiente
$f'\colon\cal{F}'=\cal{F}\times_{\cal{E}}\cal{E}'\to\allowbreak
\cal{G}'=\cal{G}\times_{\cal{E}}\cal{E}'$ es plenamente fiel.
\end{proposition}

La comprobación es inmediata; más generalmente, puede demostrarse que todo
límite proyectivo de
\ifthenelse{\boolean{orig}}{funtor}{funtores}
plenamente
\ifthenelse{\boolean{orig}}{fiel}{fieles}
(aquí,~$f$ y los funtores identidad en~$\cal{E},\cal{E}'$) es un funtor
plenamente fiel.

Obsérvese que la afirmación análoga a~4.1,
\ifthenelse{\boolean{orig}}{o}{donde}
«plenamente fiel» se sustituye por «equivalencia de categorías», es falsa ya
para~$\cal{G}=\cal{E}$. Sin embargo:

\begin{proposition}
\label{VI.4.2} Sea $f\colon\cal{F}\to\cal{G}$ un
$\cal{E}$-funtor. Las condiciones siguientes son equivalentes:
\begin{enumerate}
\item[(i)] Existe un $\cal{E}$-funtor $g\colon\cal{G}\to\cal{F}$ y
$\cal{E}$-isomorfismos
$$
gf\isomto\id_{\cal{F}},\quad fg\isomto\id_{\cal{G}}
\ifthenelse{\boolean{orig}}{}{.}
$$
\item[(ii)] Para toda categoría $\cal{E}'$ sobre~$\cal{E}$, el funtor
$$
f'=f\times_{\cal{E}}\cal{E}'\colon\cal{F}'=\cal{F}
\times_{\cal{E}}\cal{E}'\to\cal{G}'=\cal{G}\times_{\cal{E}}\cal{E}'
$$
es una equivalencia de categorías.
\item[(iii)]
\marginpar{160}
$f$ es una equivalencia de categorías y, para todo $S\in\Ob(\cal{E})$, el
funtor $f_S\colon\cal{F}_S\to\cal{G}_S$ inducido por~$f$ es una equivalencia
de categorías.
\item[(iii~bis)] $f$ es plenamente fiel y, para todo $S\in\Ob(\cal{E})$ y
todo
\ifthenelse{\boolean{orig}}
{$\eta\in\Ob\cal{G}_S$,}
{$\eta\in\Ob(\cal{G}_S)$,}
existen un $\xi\in\Ob(\cal{F}_S)$ y un $S$-isomorfismo
\ifthenelse{\boolean{orig}}
{$u\colon f((\xi)\to\eta$.}
{$u\colon f(\xi)\to\eta$.}
\end{enumerate}
\end{proposition}

\subsubsection*{Demostración} Evidentemente, (i) implica que~$f$ es una
equivalencia de categorías (noción que se define mediante la misma condición,
pero sin exigir que los isomorfismos de funtores sean $\cal{E}$-morfismos).
Por otra parte, de las funtorialidades del número anterior se deduce que la
condición~(i) se conserva después del cambio de base
$\cal{E}'\to\cal{E}$. Por consiguiente, (i)$\To$(ii). Evidentemente,
(ii)$\To$(iii), pues basta tomar $\cal{E}'=\cal{E}$ y
$\cal{E}'=\{S\}$. Es aún más trivial que (iii)$\To$(iii~bis); queda por
probar que (iii~bis)$\To$(i). Para ello, elijamos, para todo
$\eta\in\Ob(\cal{G})$, un $g(\eta)\in\Ob(\cal{F})$ y un isomorfismo
$u(\eta)\colon f(g(\eta))\to\eta$ tal que $q(u(\eta))=\id_S$, donde
$S=q(\eta)$. Esto es posible gracias a la segunda condición de (iii~bis).
Como es bien sabido e inmediato, el hecho de que~$f$ sea plenamente fiel
implica que~$g$ puede considerarse, de manera única, como un funtor de
$\cal{G}$ en~$\cal{F}$, de modo que los~$u(\eta)$ definan una transformación
natural (y, por tanto, un isomorfismo)
$u\colon fg\isomto\id_{\cal{G}}$. Además, por construcción,~$g$ es un
$\cal{E}$-funtor y~$u$ un $\cal{E}$-homomorfismo. A los datos anteriores
corresponde entonces un isomorfismo funtorial
$v\colon gf\to\id_{\cal{F}}$, definido por la condición $f*v=u*f$, y se
comprueba inmediatamente que también es un $\cal{E}$-morfismo, cqfd.

\begin{definition}
\label{VI.4.3} Si se cumplen las condiciones anteriores, se dice que~$f$ es
una equivalencia de categorías sobre~$\cal{E}$, o una
$\cal{E}$-equivalencia.
\end{definition}

\begin{corollaire}
\label{VI.4.4} Supongamos que el funtor de proyección
$p\colon\cal{F}\to\cal{E}$ es un funtor transportable, es decir, que para
todo isomorfismo $\alpha\colon T\to S$ en~$\cal{E}$ y todo objeto~$\xi$
de~$\cal{F}_T$, existe un isomorfismo~$u$ en~$\cal{F}$ de fuente~$\xi$ tal
que $p(u)=\alpha$. Entonces todo $\cal{E}$-funtor
$f\colon\cal{F}\to\cal{G}$ que sea una equivalencia de categorías es una
$\cal{E}$-equivalencia.
\end{corollaire}

Se deduce
\marginpar{161}
del criterio (iii~bis).

\begin{corollaire}
\label{VI.4.5} Sea $f\colon\cal{F}\to\cal{G}$ una
$\cal{E}$-equivalencia. Entonces, para toda categoría~$\cal{H}$
sobre~$\cal{E}$, los funtores correspondientes:
\begin{align*}
\SheafHom_{\cal{E}/-}(\cal{G},\cal{H})&\to
\SheafHom_{\cal{E}/-}(\cal{F},\cal{H})\\
\SheafHom_{\cal{E}/-}(\cal{H},\cal{F})
&\to\SheafHom_{\cal{E}/-}(\cal{H},\cal{G})
\end{align*}
(\cf \No \Ref{VI.2}) son equivalencias de categorías.
\end{corollaire}

Esto se deduce del criterio~(i) mediante el razonamiento habitual.
